% Luc Destombe --- licence CC BY-NC-SA 4.0
% https://creativecommons.org/licenses/by-nc-sa/4.0/deed.fr
% Fichier autonome : compiler avec pdflatex, deux fois (signets, nombre de pages).
\def\lyceeversion{1.0}
\documentclass[11pt,a4paper]{article}
\makeatletter
\newif\ifeleve
\elevefalse

\newif\iflycee@cm
\newif\iflycee@palatino

\RequirePackage[utf8]{inputenc}
\RequirePackage[T1]{fontenc}
\RequirePackage[french]{babel}
\RequirePackage{etoolbox}

\newcommand{\ShorthandsOff}{\@ifundefined{shorthandoff}{}{\shorthandoff{;:!?}}}
\newcommand{\ShorthandsOn}{\@ifundefined{shorthandon}{}{\shorthandon{;:!?}}}
\ShorthandsOff
\AtBeginDocument{\ShorthandsOn}
\AtBeginEnvironment{tikzpicture}{\ShorthandsOff}

\RequirePackage{geometry}
\RequirePackage{fancyhdr}
\RequirePackage{lastpage}
\RequirePackage{multicol}
\RequirePackage{array}
\RequirePackage{enumitem}
\RequirePackage{xcolor}
\RequirePackage{needspace}

\RequirePackage{amsmath,amssymb}
\iflycee@cm\else\RequirePackage{fourier}\fi
\DeclareMathAlphabet{\mathcal}{OMS}{cmsy}{m}{n}
\renewcommand{\ttdefault}{lmtt}
\begingroup\catcode`\_=\active \catcode`\^=\active
\global\let_\sb \global\let^\sp
\endgroup
\AtBeginDocument{\catcode`\_=12 \mathcode`\_="8000
  \catcode`\^=12 \mathcode`\^="8000 }
\newlength{\retraitcalculs}
\setlength{\retraitcalculs}{2em}

\RequirePackage{tikz}
\usetikzlibrary{arrows.meta,calc,babel,shapes.misc}
\RequirePackage[most]{tcolorbox}
\tcbuselibrary{skins,breakable}

\definecolor{coulDef}{HTML}{1F4E79}
\definecolor{coulProp}{HTML}{7030A0}
\definecolor{coulRem}{HTML}{C55A11}
\definecolor{coulEx}{HTML}{2E7D32}
\definecolor{coulExo}{HTML}{B03A2E}
\definecolor{coulPart}{HTML}{1F4E79}
\definecolor{coulMeth}{HTML}{1B6E4F}
\definecolor{coulCourbe}{HTML}{0B5ED7}
\definecolor{coulCourbeB}{HTML}{C00000}
\definecolor{coulCourbeC}{HTML}{2E7D32}
\definecolor{coulGrille}{HTML}{8C8C8C}
\definecolor{coulRep}{HTML}{1B6E4F}
\definecolor{coulBar}{HTML}{2E7D32}

\newcommand{\R}{\mathbb{R}}
\newcommand{\Rs}{\mathbb{R}^{*}}
\newcommand{\N}{\mathbb{N}}
\newcommand{\Ns}{\mathbb{N}^{*}}
\newcommand{\Z}{\mathbb{Z}}
\newcommand{\Q}{\mathbb{Q}}
\renewcommand{\le}{\leqslant}
\renewcommand{\ge}{\geqslant}
\renewcommand{\approx}{\simeq}
\newcommand{\vect}[1]{\overrightarrow{#1}}
\newcommand{\norme}[1]{\left\lVert #1 \right\rVert}
\newcommand{\intervalle}[2]{\left[#1\,;#2\right]}
\RequirePackage{cancel}
\renewcommand{\CancelColor}{\color{coulExo}}
\newcommand{\fois}[1]{\times{\color{coulExo}#1}}
\newcommand{\barre}[1]{\cancel{#1}}
\newlength{\lycee@eqh}\newlength{\lycee@eqd}
\newsavebox{\lycee@eqa}\newsavebox{\lycee@eqb}\newsavebox{\lycee@eqc}
\newcommand{\lycee@eqcol}[1]{\hbox{$\renewcommand{\arraystretch}{1.35}%
  \begin{array}[t]{@{}>{\displaystyle}r@{}>{\displaystyle{}}l@{}}#1\end{array}$}}
\newcommand{\lycee@eqmesure}[1]{%
  \ifdim\ht#1>\lycee@eqh\setlength{\lycee@eqh}{\ht#1}\fi
  \ifdim\dp#1>\lycee@eqd\setlength{\lycee@eqd}{\dp#1}\fi}
\newcommand{\lycee@eqtrait}{\hspace{1.2em}ou\hspace{1.2em}}
\newcommand{\lycee@eqiff}{$\iff$\hspace{0.8em}}
\newcommand{\eqproduit}[2]{\par\smallskip\noindent C'est une équation produit :
  $A\times B=0 \iff A=0$ ou $B=0$.\par\nopagebreak\smallskip
  \begingroup\setlength{\lycee@eqh}{0pt}\setlength{\lycee@eqd}{0pt}%
  \setbox\lycee@eqa=\lycee@eqcol{#1}\setbox\lycee@eqb=\lycee@eqcol{#2}%
  \lycee@eqmesure\lycee@eqa\lycee@eqmesure\lycee@eqb%
  \noindent\hspace*{\retraitcalculs}\lycee@eqiff\usebox\lycee@eqa\lycee@eqtrait\usebox\lycee@eqb\endgroup\par\smallskip}
\newcommand{\eqproduittrois}[3]{\par\smallskip\noindent C'est une équation produit :
  $A\times B\times C=0 \iff A=0$ ou $B=0$ ou $C=0$.\par\nopagebreak\smallskip
  \begingroup\setlength{\lycee@eqh}{0pt}\setlength{\lycee@eqd}{0pt}%
  \setbox\lycee@eqa=\lycee@eqcol{#1}\setbox\lycee@eqb=\lycee@eqcol{#2}\setbox\lycee@eqc=\lycee@eqcol{#3}%
  \lycee@eqmesure\lycee@eqa\lycee@eqmesure\lycee@eqb\lycee@eqmesure\lycee@eqc%
  \noindent\hspace*{\retraitcalculs}\lycee@eqiff\usebox\lycee@eqa\lycee@eqtrait\usebox\lycee@eqb\lycee@eqtrait\usebox\lycee@eqc\endgroup
  \par\smallskip}

\tikzset{grille/.style={coulGrille,line width=0.4pt}}
\newcommand{\quadrillage}[4]{\draw[grille,step=1] (#1,#2) grid (#3,#4);}
\tikzset{
  croix/.style={cross out,draw,line width=0.6pt,inner sep=0pt,minimum size=4pt},
  croixdroite/.style={croix,rotate=45,line width=0.8pt},
}
\newcommand{\pt}[3]{\path (#1) node[croix]{} node[#2,font=\small,inner sep=2.5pt]{$#3$};}

\newcommand{\lycee@repere}[4]{%
  \draw[grille,step=1] (#1,#2) grid (#3,#4);
  \draw[-{Stealth[length=2mm]},thick] (#1,0)--({#3+0.4},0) node[below left=-1pt]{$x$};
  \draw[-{Stealth[length=2mm]},thick] (0,#2)--(0,{#4+0.4}) node[below left=-1pt]{$y$};
  \node[below left,font=\scriptsize,inner sep=1.5pt] at (0,0) {$0$};}
\newcommand{\lycee@gradx}[1]{\foreach \k/\l in {#1}{%
  \draw (\k,0.07)--(\k,-0.07) node[below,font=\scriptsize,inner sep=1.5pt]{$\l$};}}
\newcommand{\lycee@grady}[1]{\foreach \k/\l in {#1}{%
  \draw (0.07,\k)--(-0.07,\k) node[left,font=\scriptsize,inner sep=1.5pt]{$\l$};}}
\newcommand{\lycee@intff}[2]{\left[#1\,;#2\right]}
\newcommand{\lycee@intfo}[2]{\left[#1\,;#2\right[}
\newcommand{\lycee@intof}[2]{\left]#1\,;#2\right]}
\newcommand{\lycee@intoo}[2]{\left]#1\,;#2\right[}
\newcommand{\lycee@ens}[1]{\left\{#1\right\}}
\AtBeginDocument{%
  \providecommand{\repere}{\lycee@repere}%
  \providecommand{\gradx}{\lycee@gradx}%
  \providecommand{\grady}{\lycee@grady}%
  \providecommand{\Cf}{\mathcal{C}\sb{f}}%
  \providecommand{\Cg}{\mathcal{C}\sb{g}}%
  \providecommand{\intff}{\lycee@intff}%
  \providecommand{\intfo}{\lycee@intfo}%
  \providecommand{\intof}{\lycee@intof}%
  \providecommand{\intoo}{\lycee@intoo}%
  \providecommand{\ens}{\lycee@ens}}

\newlist{questions}{enumerate}{3}
\setlist[questions,1]{label=\arabic*\textdegree),leftmargin=*,itemsep=2pt,topsep=3pt}
\setlist[questions,2]{label=\alph*),leftmargin=*,itemsep=1pt,topsep=2pt}
\setlist[questions,3]{label=\roman*),leftmargin=*,itemsep=1pt,topsep=1pt}
\setlist[itemize]{label=\textbullet,leftmargin=*,itemsep=2pt,topsep=3pt}

\newcommand{\besoinplace}[1]{\par\penalty\@M\begingroup
  \dimen@=#1\relax\dimen@ii=\pagegoal\advance\dimen@ii-\pagetotal
  \ifdim\dimen@>\dimen@ii\ifdim\dimen@ii>\z@\vfil\fi\break\fi\endgroup}

\newcommand{\lycee@pied}{\thepage/\pageref{LastPage}}
\newcommand{\entete}[2]{%
  \pagestyle{fancy}\fancyhf{}%
  \fancyhead[L]{\footnotesize\color{coulPart}\textbf{#1}}%
  \fancyhead[R]{\footnotesize\color{coulPart}\textbf{#2}}%
  \fancyfoot[C]{\footnotesize\lycee@pied}%
  \renewcommand{\headrulewidth}{0.4pt}%
  \renewcommand{\headrule}{\hbox to\headwidth{%
    \color{coulPart!40}\leaders\hrule height \headrulewidth\hfill}}}

\RequirePackage{ccicons}
\newcommand{\lycee@licence}{{\scriptsize\color{gray}%
  \copyright\ Luc Destombe --- \ccbyncsa\ CC BY-NC-SA 4.0}}
\newif\iflycee@licence \lycee@licencetrue
\newcommand{\sanslicence}{\lycee@licencefalse}
\AtBeginDocument{\iflycee@licence\AddToHookNext{shipout/foreground}{%
  \put(0,0){\rlap{\kern\dimexpr1in+\hoffset+\oddsidemargin+\textwidth\relax
    \llap{\raisebox{-\dimexpr1in+\voffset+\topmargin+\headheight+\headsep
      +\textheight+\footskip\relax}{\lycee@licence}}}}}\fi}

\newcommand{\lycee@site}{faire-des-maths.fr}
\newcommand{\lycee@version}{{\scriptsize\color{gray}\lycee@site
  \ifdefined\lyceeversion\ --- Version~\lyceeversion\fi}}
\AtBeginDocument{\iflycee@licence\AddToHookNext{shipout/foreground}{%
  \put(0,0){\rlap{\kern\dimexpr1in+\hoffset+\oddsidemargin\relax
    \raisebox{-\dimexpr1in+\voffset+\topmargin+\headheight+\headsep
      +\textheight+\footskip\relax}{\lycee@version}}}}\fi}

\newcommand{\formulesserrees}{%
  \setlength{\abovedisplayskip}{3pt}\setlength{\belowdisplayskip}{3pt}%
  \setlength{\abovedisplayshortskip}{2pt}\setlength{\belowdisplayshortskip}{3pt}}

\setlength{\parindent}{0pt}

\RequirePackage[implicit=false,hidelinks,pdfencoding=auto,psdextra,bookmarksopen,
  bookmarksopenlevel=1,pdfcreator={},pdfproducer={}]{hyperref}
\RequirePackage{bookmark}
\hypersetup{pdfauthor={Luc Destombe},pdfkeywords={CC BY-NC-SA 4.0}}
\newcounter{lycee@signet}
\newcommand{\signet}[3][]{\stepcounter{lycee@signet}%
  \edef\lycee@dest{\if\relax\detokenize{#1}\relax
    signet.\arabic{lycee@signet}\else#1\fi}%
  \hypertarget{\lycee@dest}{}\bookmark[level=#2,dest=\lycee@dest]{#3}}
\pdfstringdefDisableCommands{%
  \def\R{R}\def\N{N}\def\Z{Z}\def\Q{Q}\def\vect#1{#1}\def\Cf{Cf}\def\Cg{Cg}\def\fois#1{×#1}%
  \def\barre#1{#1}\def\intervalle#1#2{[#1;#2]}\def\intff#1#2{[#1;#2]}%
  \def\textsuperscript#1{#1}\def\dfrac#1#2{#1/#2}\def\frac#1#2{#1/#2}%
  \def\vec#1{#1⃗}}%
\newcommand{\res}[1]{#1}
\geometry{top=2cm,bottom=1cm,left=1cm,right=1cm,headheight=14pt,footskip=0.6cm}

\RequirePackage{pgfplots}
\pgfplotsset{compat=1.18}
\RequirePackage{tkz-tab}

\newcounter{partiecours}
\renewcommand{\thepartiecours}{\Roman{partiecours}}
\newcommand{\partiecours}[1]{%
  \refstepcounter{partiecours}%
  \Needspace*{8\baselineskip}%
  \bigskip
  \noindent\begin{tcolorbox}[enhanced,colback=coulPart,colframe=coulPart,
      boxrule=0pt,arc=2pt,left=8pt,right=8pt,top=4pt,bottom=4pt,
      before skip=6pt,after skip=10pt]
    \color{white}\bfseries\large \thepartiecours{} --- #1%
    \signet[partie-\arabic{partiecours}]{1}{\thepartiecours{} --- #1}
  \end{tcolorbox}%
  \addcontentsline{toc}{section}{\thepartiecours{} --- #1}%
}

\newtcolorbox{boitecours}[3][]{%
  enhanced,breakable,
  colback=white,colframe=#2,coltitle=white,
  fonttitle=\bfseries,
  attach boxed title to top left={xshift=8pt,yshift=-\tcboxedtitleheight/2},
  boxed title style={colback=#2,arc=2pt,boxrule=0pt},
  boxrule=0.9pt,arc=2pt,
  left=8pt,right=8pt,top=8pt,bottom=6pt,
  title={#3},#1}

\newcounter{methode}
\newenvironment{definitionbox}[1][Définition]%
  {\begin{boitecours}[phantom={\signet{2}{#1}}]{coulDef}{#1}}{\end{boitecours}}
\newenvironment{proprietebox}[1][Propriété]%
  {\begin{boitecours}[phantom={\signet{2}{#1}}]{coulProp}{#1}}{\end{boitecours}}
\newenvironment{methodebox}[1][Méthode]{\stepcounter{methode}%
  \begin{boitecours}[phantom={\signet[methode-\arabic{methode}]{2}{#1}}]{coulMeth}{#1}}%
  {\end{boitecours}}

\newcommand{\etape}[1]{\par\smallskip\textbf{\color{coulMeth}#1}\par\nobreak\smallskip}
\newcommand{\enonce}[1]{\emph{#1}\par\smallskip}
\newcommand{\reponse}[1]{\par\smallskip
  {\footnotesize\itshape\color{black!55}Réponses : #1}}
\newcommand{\demo}[1]{\textbf{\color{coulMeth}Démonstration.} #1}
\newcommand{\afaire}[1]{%
  \par\smallskip
  \noindent{\small\color{coulExo}$\blacktriangleright$\ \textbf{À faire :}
    fiche d'exercices du chapitre, #1.}\par\smallskip}

\newenvironment{calculs}{\setlength{\jot}{5pt}\@fleqntrue\@mathmargin\retraitcalculs
  \setlength{\abovedisplayskip}{4pt}\setlength{\belowdisplayskip}{4pt}%
  \setlength{\abovedisplayshortskip}{4pt}\setlength{\belowdisplayshortskip}{4pt}%
  \csname alignat*\endcsname{2}}%
  {\csname endalignat*\endcsname}
\newcommand{\rem}[1]{\text{\small\itshape\color{black!60}#1}}
\newcommand{\explique}[2]{\par\smallskip\noindent
  \begin{minipage}[t]{0.57\linewidth}\raggedright #1\end{minipage}\hfill
  \begin{minipage}[t]{0.40\linewidth}\raggedright\leavevmode
    {\small\itshape\color{black!60}#2}\end{minipage}\par}
\newcommand{\cas}[1]{\par\smallskip\noindent\textbf{\color{coulExo}#1}\vadjust{\nobreak}}
\newcommand{\calc}[1]{\par\smallskip\textbf{\color{coulExo}#1}\par\nobreak}

\newtcolorbox{remarquebox}[1][Remarque]{%
  enhanced,breakable,colback=white,colframe=coulRem,
  boxrule=0pt,leftrule=2.5pt,arc=0pt,outer arc=0pt,
  left=8pt,right=6pt,top=5pt,bottom=5pt,
  before upper={\textbf{\color{coulRem}#1}\par\smallskip}}

\newtcolorbox{exemplebox}[1][Exemple]{%
  enhanced,breakable,colback=white,colframe=coulEx,
  boxrule=0pt,leftrule=2.5pt,arc=0pt,outer arc=0pt,
  left=8pt,right=6pt,top=5pt,bottom=5pt,
  before upper={\textbf{\color{coulEx}#1}\par\smallskip}}

\pgfplotsset{
  repere/.style={
    axis lines=middle,
    axis line style={-{Stealth[length=2.4mm]},thick},
    every axis x label/.style={at={(ticklabel* cs:1.0)},anchor=west},
    every axis y label/.style={at={(ticklabel* cs:1.0)},anchor=south},
    label style={font=\footnotesize},
    tick label style={font=\scriptsize},
    grid=both,
    major grid style={grille},
    minor tick num=0,
    tick align=inside,
    clip mode=individual,
  },
}
\tikzset{
  courbe/.style={coulCourbe,line width=1.1pt,smooth},
  courbeB/.style={coulCourbeB,line width=1.1pt,smooth},
  courbeC/.style={coulCourbeC,line width=1.1pt,smooth},
}

\newlist{etapes}{enumerate}{1}
\setlist[etapes]{label=\textbf{\arabic*.},leftmargin=*,itemsep=1pt,topsep=2pt}

\newlist{expressions}{enumerate}{1}
\setlist[expressions]{label={\Alph*\ $=$},leftmargin=*,itemsep=3pt,topsep=3pt}

\newcounter{exo}
\renewcommand{\theexo}{\ifnum\value{exo}<10 0\fi\arabic{exo}}
\newtcolorbox{exercice}[1][]{%
  enhanced,breakable,
  colback=white,colframe=coulExo,
  boxrule=0.9pt,arc=2pt,
  left=8pt,right=8pt,top=8pt,bottom=6pt,
  coltitle=white,fonttitle=\bfseries,
  attach boxed title to top left={xshift=8pt,yshift=-\tcboxedtitleheight/2},
  boxed title style={colback=coulExo,arc=2pt,boxrule=0pt},
  phantom={\signet{2}{Exercice \arabic{exo}}},
  title={Exercice \theexo\ifx\relax#1\relax\else{} --- #1\fi}}
\newenvironment{exo}[1][\relax]{\refstepcounter{exo}\begin{exercice}[#1]}{\end{exercice}}

  \newenvironment{correction}{\par}{\par}
  \newcommand{\autableau}{}
  \newcommand{\etapeexemples}{\etape{Exemples corrigés}}
  \newcommand{\etapetableau}[1]{\etape{#1}}
  \newenvironment{exempletableau}[1][Exemple]%
    {\begin{exemplebox}[#1]}{\end{exemplebox}}

\RequirePackage{comment}
\excludecomment{correctionavj}

\newcommand{\coursEntete}{}
\newcommand{\coursNiveau}{}
\newcommand{\coursTitre}{}
\newcommand{\coursSurTitre}{}
\newcommand{\coursSousTitre}{}

\newcommand{\enteteCours}[1]{\renewcommand{\coursEntete}{#1}}
\newcommand{\chapitrecours}[2]{\enteteCours{Chapitre #1 --- #2}}
\newcommand{\niveaucours}[1]{\renewcommand{\coursNiveau}{#1}}
\newcommand{\titrecours}[3][]{%
  \renewcommand{\coursTitre}{#2}%
  \renewcommand{\coursSurTitre}{#1}%
  \renewcommand{\coursSousTitre}{#3}}

\pagestyle{fancy}
\fancyhf{}
\fancyhead[L]{\footnotesize\color{coulPart}\textbf{\coursEntete}}
\fancyhead[R]{\footnotesize\color{coulPart}\textbf{\coursNiveau}}
\fancyfoot[C]{\footnotesize\thepage}
\renewcommand{\headrulewidth}{0.4pt}
\renewcommand{\headrule}{\hbox to\headwidth{\color{coulPart!40}\leaders\hrule height \headrulewidth\hfill}}

\setlength{\parskip}{2pt}

\newcommand{\lycee@mentioneleve}{\\[2pt]{\small\itshape Version élève}}
\AtBeginDocument{%
  \begin{center}
    \begin{tcolorbox}[enhanced,width=0.72\linewidth,colback=white,colframe=coulPart,
        boxrule=1.6pt,arc=3pt,halign=center,top=10pt,bottom=10pt]
      \ifx\coursSurTitre\empty
        {\Huge\bfseries\color{coulPart} \coursTitre}\\[4pt]
      \else
        {\Huge\bfseries\color{coulPart} \coursTitre}\\[3pt]
        {\Large\bfseries\color{coulPart} \coursSurTitre}\\[5pt]
      \fi
      {\normalsize \coursSousTitre}%
      
    \end{tcolorbox}
  \end{center}%
}

\RequirePackage{tikz}
\RequirePackage{xcolor}

\usetikzlibrary{decorations.text}
\definecolor{lycee@metal}{HTML}{1F4E79}
\definecolor{lycee@manche}{HTML}{8A5A2B}
\definecolor{lycee@etincelle}{HTML}{F28C28}
\newcommand{\lycee@pluriel}[2]{#1 #2\ifnum#1>1 s\fi}
\newcommand{\lycee@sceau}{%
  \edef\lycee@bas{\lycee@pluriel{\lyceerelectures}{relecture} et
    \lycee@pluriel{\lyceerodages}{rodage}}%
  \begin{tikzpicture}[line join=round]
    \fill[coulPart!10,draw=coulPart,line width=.8pt] (0:2.2)
      \foreach \i in {1,...,71}{ -- (\i*5:{isodd(\i)?2.1:2.2}) } -- cycle;
    \draw[coulPart,line width=1.2pt] (0,0) circle (1.93);
    \draw[coulPart,line width=.4pt] (0,0) circle (1.85);
    \draw[coulPart,line width=.4pt] (0,0) circle (1.2);
    \path[decorate,decoration={text along path,text align=center,
        text={|\fontsize{9.5}{11}\selectfont\bfseries\color{coulPart}|TEST{\'E} AVEC MES {\'E}L{\`E}VES}}]
      (180:1.45) arc (180:0:1.45);
    \path[decorate,decoration={text along path,text align=center,
        text/.expanded={|\noexpand\fontsize{8.5}{10}\noexpand\selectfont\noexpand\color{coulPart}|\lycee@bas}}]
      (195:1.68) arc (195:345:1.68);
    \begin{scope}[x=.25mm,y=-.25mm]
      \begin{scope}[shift={(-32,-40)}]%
        \fill[lycee@metal] (23,31) rectangle (41,49);
        \draw[lycee@etincelle,line width=.5pt,line cap=round]
          (20,30)--(15,27) (19,34)--(14,35) (44,30)--(49,27) (45,34)--(50,35);
        \fill[lycee@manche,rounded corners=.3mm] (38,24.6) rectangle (60,27.6);
        \fill[lycee@metal,rounded corners=.4mm] (25,20) rectangle (39,31);
        \fill[lycee@etincelle,rounded corners=.2mm] (25,20) rectangle (28,31);
      \end{scope}
    \end{scope}
    \ifdefined\lyceesceauplusieurs\def\lycee@annee{Ann{\'e}es~:}\else
      \def\lycee@annee{Ann{\'e}e}\fi
    \path[decorate,decoration={text along path,text align=center,
        text/.expanded={|\noexpand\fontsize{8}{10}\noexpand\selectfont\noexpand\color{coulPart}|%
          \unexpanded\expandafter{\lycee@annee} \unexpanded\expandafter{\lyceesceauannees}}}]
      (95:1.08) arc (95:445:1.08);
  \end{tikzpicture}}
\AtEndDocument{\ifdefined\lyceesceauannees
  \par\besoinplace{5cm}\vspace{1.5em}%
  \begin{center}\lycee@sceau\end{center}\fi}
\makeatother

\chapitrecours{6}{Dérivation, convexité, continuité}
\niveaucours{T\textsuperscript{le} spécialité}
\titrecours{DÉRIVATION, CONVEXITÉ, CONTINUITÉ}{Cours --- Terminale spécialité}

\begin{document}

\newcommand{\e}{\mathrm{e}}

\newtcblisting{python}{%
  listing only,enhanced,breakable,
  colback=black!3,colframe=black!35,boxrule=0.5pt,arc=1pt,
  left=6pt,right=6pt,top=2pt,bottom=2pt,
  listing options={language=Python,basicstyle=\ttfamily\small,
    keywordstyle=\color{coulMeth}\bfseries,columns=fullflexible,
    keepspaces=true,showstringspaces=false,aboveskip=0pt,belowskip=0pt}}

\begin{remarquebox}[Durée du chapitre]
À titre indicatif, ce chapitre est prévu pour un volume de $14$~h d'enseignement
($7$ séances de $2$~h). Une partie importante des exercices est à faire en autonomie.
\end{remarquebox}

\partiecours{Rappels sur la dérivation}

En première, on a appris à dériver les fonctions usuelles et leurs sommes, produits et
quotients, puis à s'en servir pour les tangentes et les variations. Ces outils servent
dans tout le chapitre : on les rassemble ici.

\begin{proprietebox}[Propriété 1 --- Dérivées des fonctions usuelles et opérations]
\begin{minipage}[t]{0.52\linewidth}\centering
\renewcommand{\arraystretch}{2.4}\setlength{\tabcolsep}{8pt}
\begin{tabular}[t]{|c|c|c|}
\hline
$f(x)$ & $f'(x)$ & dérivable sur\\ \hline
$k$ (constante) & $0$ & $\R$\\ \hline
$x^{n}$ ($n\in\Ns$) & $nx^{n-1}$ & $\R$\\ \hline
$\dfrac{1}{x}$ & $-\dfrac{1}{x^{2}}$ & $\intoo{-\infty}{0}$, $\intoo{0}{+\infty}$\\ \hline
$\sqrt{x}$ & $\dfrac{1}{2\sqrt{x}}$ & $\intoo{0}{+\infty}$\\ \hline
$\e^{x}$ & $\e^{x}$ & $\R$\\ \hline
$\e^{ax+b}$ & $a\e^{ax+b}$ & $\R$\\ \hline
\end{tabular}
\end{minipage}\hfill
\begin{minipage}[t]{0.44\linewidth}\centering
\renewcommand{\arraystretch}{2.4}\setlength{\tabcolsep}{12pt}
\begin{tabular}[t]{|c|c|}
\hline
fonction & dérivée\\ \hline
$u+v$ & $u'+v'$\\ \hline
$ku$ ($k\in\R$) & $ku'$\\ \hline
$uv$ & $u'v+uv'$\\ \hline
$\dfrac{1}{v}$ & $-\dfrac{v'}{v^{2}}$\\ \hline
$\dfrac{u}{v}$ & $\dfrac{u'v-uv'}{v^{2}}$\\ \hline
\end{tabular}
\end{minipage}

\smallskip
Dans le second tableau, $u$ et $v$ sont dérivables sur un intervalle $I$, et $v$ ne
s'annule pas sur $I$ pour $\frac{1}{v}$ et $\frac{u}{v}$.
\end{proprietebox}

\begin{proprietebox}[Propriété 2 --- Tangente et sens de variation]
Soit $f$ une fonction dérivable sur un intervalle $I$ et $a\in I$.
\begin{itemize}
  \item La tangente $(T_{a})$ à $\Cf$ au point d'abscisse $a$ a pour équation
        $y=f'(a)(x-a)+f(a)$.
  \item Si $f'(x)>0$ pour tout $x\in I$, sauf peut-être en des valeurs isolées où
        $f'(x)=0$, alors $f$ est strictement croissante sur $I$ ; si $f'(x)<0$ dans les
        mêmes conditions, $f$ est strictement décroissante sur $I$.
\end{itemize}
\end{proprietebox}

\begin{methodebox}[Méthode 1 --- Calculer une dérivée]
\etapeexemples
Calculer la dérivée de chaque fonction, et l'écrire sous forme factorisée.
\begin{questions}[label=\alph*)]
  \item La fonction $f$ définie sur $\R$ par $f(x)=(2x-3)\e^{x}$.
  \item La fonction $g$ définie sur $\intoo{2}{+\infty}$ par $g(x)=\dfrac{3x+1}{x-2}$.
  \item La fonction $h$ définie sur $\R$ par $h(x)=x^{2}\e^{-x}$.
\end{questions}
\begin{correction}
\cas{a)}
\explique{$f$ est dérivable sur $\R$, comme produit de fonctions dérivables sur
$\R$. Pour tout $x\in\R$ :}{forme $uv$ avec $u(x)=2x-3$ et $v(x)=\e^{x}$}
\begin{calculs}
f'(x) &= 2\times\e^{x}+(2x-3)\times\e^{x} &\qquad& \rem{$u'v+uv'$, avec $u'(x)=2$ et $v'(x)=\e^{x}$}\\
f'(x) &= (2+2x-3)\e^{x} &\qquad& \rem{on factorise par $\e^{x}$}\\
f'(x) &= \res{(2x-1)\e^{x}} &&
\end{calculs}
\cas{b)}
\explique{$g$ est dérivable sur $\intoo{2}{+\infty}$, comme quotient de fonctions
dérivables dont le dénominateur ne s'y annule pas. Pour tout $x>2$ :}{forme
$\frac{u}{v}$ avec $u(x)=3x+1$ et $v(x)=x-2$}
\begin{calculs}
g'(x) &= \frac{3(x-2)-(3x+1)\times 1}{(x-2)^{2}} &\qquad& \rem{$\frac{u'v-uv'}{v^{2}}$, avec $u'(x)=3$ et $v'(x)=1$}\\
g'(x) &= \frac{3x-6-3x-1}{(x-2)^{2}} &\qquad& \rem{on développe le numérateur}\\
g'(x) &= \res{\frac{-7}{(x-2)^{2}}} &&
\end{calculs}
\cas{c)}
\explique{$h$ est dérivable sur $\R$, comme produit de fonctions dérivables sur
$\R$. Pour tout $x\in\R$ :}{$u(x)=x^{2}$ et $v(x)=\e^{-x}$, de dérivée $-\e^{-x}$}
\begin{calculs}
h'(x) &= 2x\e^{-x}+x^{2}\times\left(-\e^{-x}\right) &\qquad& \rem{$u'v+uv'$}\\
h'(x) &= \left(2x-x^{2}\right)\e^{-x} &\qquad& \rem{on factorise par $\e^{-x}$}\\
h'(x) &= \res{x(2-x)\e^{-x}} &\qquad& \rem{puis par $x$ : le signe se lit mieux}
\end{calculs}
\end{correction}

\etape{À vous de jouer}
Calculer la dérivée de chaque fonction, et l'écrire sous forme factorisée.
\begin{questions}[label=\alph*)]
  \item La fonction $f$ définie sur $\R$ par $f(x)=(3x+2)\e^{x}$.
  \item La fonction $g$ définie sur $\intoo{-1}{+\infty}$ par $g(x)=\dfrac{2x-5}{x+1}$.
  \item La fonction $h$ définie sur $\R$ par $h(x)=x^{2}\e^{-2x}$.
\end{questions}
\reponse{a) $f'(x)=(3x+5)\e^{x}$ ;\; b)~$g'(x)=\dfrac{7}{(x+1)^{2}}$ ;\;
c)~$h'(x)=2x(1-x)\e^{-2x}$.}
\end{methodebox}

\begin{methodebox}[Méthode 2 --- Déterminer une équation d'une tangente]
\etapeexemples
\begin{questions}[label=\alph*)]
  \item Soit la fonction $f$ définie sur $\R$ par $f(x)=(2x-3)\e^{x}$ ; on a
        $f'(x)=(2x-1)\e^{x}$. Déterminer une équation de la tangente $(T_{0})$ à
        $\Cf$ au point d'abscisse $0$.
  \item Soit la fonction $g$ définie sur $\intoo{2}{+\infty}$ par
        $g(x)=\dfrac{3x+1}{x-2}$ ; on a $g'(x)=\dfrac{-7}{(x-2)^{2}}$. Déterminer une
        équation de la tangente $(T_{3})$ à $\Cg$ au point d'abscisse $3$.
\end{questions}
\begin{correction}
\cas{a)}
\begin{calculs}
f(0) &= (2\times 0-3)\e^{0} = -3 &\qquad& \rem{$\e^{0}=1$}\\
f'(0) &= (2\times 0-1)\e^{0} = -1 &\qquad& \rem{le coefficient directeur}
\end{calculs}
\explique{$(T_{0})$ a pour équation $y=f'(0)(x-0)+f(0)$ :}{la formule de la
propriété~2}
\begin{calculs}
y &= -1\times(x-0)+(-3) &\qquad& \rem{on remplace}\\
y &= \res{-x-3} &&
\end{calculs}
\cas{b)}
\begin{calculs}
g(3) &= \frac{3\times 3+1}{3-2} = 10 &&\\
g'(3) &= \frac{-7}{(3-2)^{2}} = -7 &&
\end{calculs}
\explique{$(T_{3})$ a pour équation $y=g'(3)(x-3)+g(3)$ :}{}
\begin{calculs}
y &= -7(x-3)+10 &\qquad& \rem{on remplace}\\
y &= -7x+21+10 &\qquad& \rem{on développe}\\
y &= \res{-7x+31} &&
\end{calculs}
\end{correction}

\etape{À vous de jouer}
\begin{questions}[label=\alph*)]
  \item Soit la fonction $f$ définie sur $\R$ par $f(x)=(3x+2)\e^{x}$ ; on a
        $f'(x)=(3x+5)\e^{x}$. Déterminer une équation de la tangente $(T_{0})$ à
        $\Cf$ au point d'abscisse $0$.
  \item Soit la fonction $g$ définie sur $\intoo{-1}{+\infty}$ par
        $g(x)=\dfrac{2x-5}{x+1}$ ; on a $g'(x)=\dfrac{7}{(x+1)^{2}}$. Déterminer une
        équation de la tangente $(T_{1})$ à $\Cg$ au point d'abscisse $1$.
\end{questions}
\reponse{a) $(T_{0})\colon y=5x+2$ ;\; b)~$(T_{1})\colon y=\frac{7}{4}x-\frac{13}{4}$.}
\end{methodebox}

\afaire{exercices 1 à 4, puis l\'entraînement au calcul de dérivées, exercices 43 à 45}

\partiecours{Fonctions composées}

Pour calculer $\sqrt{x-3}$, on calcule d'abord $x-3$, puis on prend la racine carrée du
résultat : $x\longmapsto x-3\longmapsto\sqrt{x-3}$. Deux fonctions s'enchaînent ; on a
déjà rencontré cet enchaînement au chapitre~5, pour les limites (« on pose $X=x-3$ »).

\begin{definitionbox}[Définition 1 --- Composée de deux fonctions]
Soit $u$ une fonction définie sur un intervalle $I$, et $v$ une fonction définie sur un
intervalle $J$ qui contient toutes les valeurs $u(x)$ pour $x\in I$. La \textbf{composée}
de $u$ suivie de $v$ est la fonction notée $v\circ u$ (lire « $v$ rond $u$ »), définie
sur $I$ par
\begin{center}
$v\circ u(x)=v\bigl(u(x)\bigr)$.
\end{center}
On applique d'abord $u$, puis $v$ : $x\longmapsto u(x)\longmapsto v\bigl(u(x)\bigr)$.
L'ordre compte : en général, $v\circ u$ et $u\circ v$ sont deux fonctions différentes.
\end{definitionbox}

\begin{methodebox}[Méthode 3 --- Composer et décomposer des fonctions]
\etapeexemples
\begin{questions}[label=\alph*)]
  \item Soit les fonctions $u$ et $v$ définies sur $\R$ par $u(x)=3x-1$ et $v(x)=x^{2}$.
        Exprimer $v\circ u(x)$ et $u\circ v(x)$ en fonction de $x$.
  \item Soit la fonction $f$ définie sur $\R$ par $f(x)=\e^{-x^{2}+4}$. Écrire $f$ comme
        la composée de deux fonctions $u$ et $v$ à préciser.
\end{questions}
\begin{correction}
\cas{a)}
\begin{calculs}
v\circ u(x) &= v\bigl(u(x)\bigr) &\qquad& \rem{d'abord $u$, puis $v$}\\
v\circ u(x) &= v(3x-1) &\qquad& \rem{$v$ élève au carré ce qu'on lui donne}\\
v\circ u(x) &= \res{(3x-1)^{2}} &&\\
u\circ v(x) &= u\bigl(v(x)\bigr) &\qquad& \rem{d'abord $v$, puis $u$}\\
u\circ v(x) &= u\left(x^{2}\right) &\qquad& \rem{$u$ multiplie par $3$ et retire $1$}\\
u\circ v(x) &= \res{3x^{2}-1} &&
\end{calculs}
\explique{Les deux fonctions sont différentes : en $x=1$, on obtient $4$ et $2$.}{l'ordre
compte}
\cas{b)}
\explique{Pour calculer $f(x)$, on calcule d'abord $-x^{2}+4$, puis on prend
l'exponentielle du résultat : $f=v\circ u$ avec \res{$u(x)=-x^{2}+4$ et
$v(x)=\e^{x}$}.}{$x\longmapsto -x^{2}+4\longmapsto \e^{-x^{2}+4}$}
\end{correction}

\etape{À vous de jouer}
\begin{questions}[label=\alph*)]
  \item Soit les fonctions $u$ et $v$ définies sur $\R$ par $u(x)=2x+5$ et $v(x)=x^{3}$.
        Exprimer $v\circ u(x)$ et $u\circ v(x)$ en fonction de $x$.
  \item Soit la fonction $g$ définie sur $\R$ par $g(x)=\sqrt{x^{2}+9}$. Écrire $g$
        comme la composée de deux fonctions $u$ et $v$ à préciser.
\end{questions}
\reponse{a) $v\circ u(x)=(2x+5)^{3}$ et $u\circ v(x)=2x^{3}+5$ ;\;
b)~$g=v\circ u$ avec $u(x)=x^{2}+9$ et $v(x)=\sqrt{x}$.}
\end{methodebox}

\begin{proprietebox}[Théorème 1 --- Dérivée d'une fonction composée (admis)]
Soit $u$ une fonction dérivable sur un intervalle $I$, et $v$ une fonction dérivable sur
un intervalle $J$ qui contient toutes les valeurs $u(x)$ pour $x\in I$. Alors $v\circ u$
est dérivable sur $I$ et
\begin{center}
$(v\circ u)'=\left(v'\circ u\right)\times u'$, \quad c'est-à-dire \quad
$\Bigl(v\bigl(u(x)\bigr)\Bigr)'=v'\bigl(u(x)\bigr)\times u'(x)$.
\end{center}
On dérive la fonction « extérieure » $v$ en gardant $u(x)$ à l'intérieur, puis on
multiplie par la dérivée de l'intérieur.
\end{proprietebox}

\begin{proprietebox}[Propriété 3 --- Cas particuliers à connaître]
Soit $u$ une fonction dérivable sur un intervalle $I$.
\begin{center}
\renewcommand{\arraystretch}{2.4}\setlength{\tabcolsep}{14pt}
\begin{tabular}{|c|c|c|}
\hline
fonction & dérivée & condition\\ \hline
$\e^{u}$ & $u'\e^{u}$ & \\ \hline
$u^{n}$ ($n\in\Ns$) & $nu'u^{n-1}$ & \\ \hline
$\sqrt{u}$ & $\dfrac{u'}{2\sqrt{u}}$ & $u(x)>0$ pour tout $x\in I$\\ \hline
$x\longmapsto f(ax+b)$ & $x\longmapsto af'(ax+b)$ & \\ \hline
\end{tabular}
\end{center}
\demo{pour $\sqrt{u}$, on applique le théorème~1 avec $v(X)=\sqrt{X}$, de dérivée
$v'(X)=\frac{1}{2\sqrt{X}}$ sur $\intoo{0}{+\infty}$.}
\begin{calculs}
\left(\sqrt{u(x)}\right)' &= v'\bigl(u(x)\bigr)\times u'(x) &\qquad& \rem{théorème~1}\\
\left(\sqrt{u(x)}\right)' &= \frac{1}{2\sqrt{u(x)}}\times u'(x) &\qquad& \rem{on remplace $X$ par $u(x)$}\\
\left(\sqrt{u(x)}\right)' &= \frac{u'(x)}{2\sqrt{u(x)}} &&
\end{calculs}
Les autres formules se démontrent de même, avec $v(X)=\e^{X}$, $v(X)=X^{n}$ et
$u(x)=ax+b$.
\end{proprietebox}

\begin{methodebox}[Méthode 4 --- Dériver une fonction composée]
\etapeexemples
Calculer la dérivée de chaque fonction.
\begin{questions}[label=\alph*)]
  \item La fonction $f$ définie sur $\R$ par $f(x)=\e^{x^{2}-3x}$.
  \item La fonction $g$ définie sur $\R$ par $g(x)=(4x+1)^{5}$.
  \item La fonction $h$ définie sur $\R$ par $h(x)=\sqrt{x^{2}+2x+5}$.
\end{questions}
\begin{correction}
\cas{a)}
\explique{On pose $u(x)=x^{2}-3x$ : $u$ est dérivable sur $\R$ et $u'(x)=2x-3$. Donc
$f=\e^{u}$ est dérivable sur $\R$ et, pour tout $x\in\R$ :}{on repère la forme
$\e^{u}$}
\begin{calculs}
f'(x) &= u'(x)\e^{u(x)} &\qquad& \rem{$\left(\e^{u}\right)'=u'\e^{u}$}\\
f'(x) &= \res{(2x-3)\e^{x^{2}-3x}} &&
\end{calculs}
\cas{b)}
\explique{On pose $u(x)=4x+1$ : $u'(x)=4$. Donc $g=u^{5}$ est dérivable sur $\R$ et,
pour tout $x\in\R$ :}{forme $u^{n}$ avec $n=5$}
\begin{calculs}
g'(x) &= 5u'(x)\bigl(u(x)\bigr)^{4} &\qquad& \rem{$\left(u^{n}\right)'=nu'u^{n-1}$}\\
g'(x) &= 5\times 4\times(4x+1)^{4} &&\\
g'(x) &= \res{20(4x+1)^{4}} &&
\end{calculs}
\cas{c)}
\explique{On pose $u(x)=x^{2}+2x+5$. Pour tout $x\in\R$, $u(x)=(x+1)^{2}+4>0$ ; $u$ est
dérivable sur $\R$ et $u'(x)=2x+2$. Donc $h=\sqrt{u}$ est dérivable sur $\R$ et, pour
tout $x\in\R$ :}{on vérifie $u>0$ avant d'écrire $\sqrt{u}$}
\begin{calculs}
h'(x) &= \frac{u'(x)}{2\sqrt{u(x)}} &\qquad& \rem{$\left(\sqrt{u}\right)'=\frac{u'}{2\sqrt{u}}$}\\
h'(x) &= \frac{2x+2}{2\sqrt{x^{2}+2x+5}} &&\\
h'(x) &= \frac{\barre{2}(x+1)}{\barre{2}\sqrt{x^{2}+2x+5}} &\qquad& \rem{on simplifie par $2$}\\
h'(x) &= \res{\frac{x+1}{\sqrt{x^{2}+2x+5}}} &&
\end{calculs}
\end{correction}

\etape{À vous de jouer}
Calculer la dérivée de chaque fonction.
\begin{questions}[label=\alph*)]
  \item La fonction $f$ définie sur $\R$ par $f(x)=\e^{-x^{2}+5x}$.
  \item La fonction $g$ définie sur $\R$ par $g(x)=(3x-2)^{4}$.
  \item La fonction $h$ définie sur $\R$ par $h(x)=\sqrt{x^{2}-4x+13}$.
\end{questions}
\reponse{a) $f'(x)=(-2x+5)\e^{-x^{2}+5x}$ ;\; b)~$g'(x)=12(3x-2)^{3}$ ;\;
c)~$h'(x)=\dfrac{x-2}{\sqrt{x^{2}-4x+13}}$.}
\end{methodebox}

\afaire{exercices 5 à 9, puis l\'entraînement au calcul de dérivées, exercices 46 à 48}

\partiecours{Étude d'une fonction composée}

Les sujets de bac enchaînent les trois outils : limites (chapitre~5), dérivée d'une
composée, signe de la dérivée. Le facteur $\e^{u(x)}$ est toujours strictement positif :
le signe de la dérivée est celui de l'autre facteur.

\begin{methodebox}[Méthode 5 --- Étudier une fonction composée]
\etapeexemples
Soit la fonction $f$ définie sur $\R$ par $f(x)=\e^{-x^{2}+2x}$.
\begin{questions}[label=\alph*)]
  \item Déterminer les limites de $f$ en $-\infty$ et en $+\infty$.
  \item Calculer $f'(x)$ pour tout $x\in\R$.
  \item Dresser le tableau de variations de $f$.
\end{questions}
\begin{correction}
\cas{a)}
\explique{En $+\infty$ : pour tout $x\in\R$, $-x^{2}+2x=x(2-x)$. Or
$\lim\limits_{x\to+\infty}x=+\infty$ et $\lim\limits_{x\to+\infty}(2-x)=-\infty$, donc
par produit, $\lim\limits_{x\to+\infty}\left(-x^{2}+2x\right)=-\infty$.}{forme
« $-\infty+\infty$ » : on factorise par $x$}
\explique{On pose $X=-x^{2}+2x$. Comme $\lim\limits_{X\to-\infty}\e^{X}=0$, par
composition, \res{$\lim\limits_{x\to+\infty}f(x)=0$}.}{puis l'exponentielle}
\explique{En $-\infty$ : $\lim\limits_{x\to-\infty}x=-\infty$ et
$\lim\limits_{x\to-\infty}(2-x)=+\infty$, donc par produit,
$\lim\limits_{x\to-\infty}\left(-x^{2}+2x\right)=-\infty$ ; par composition,
\res{$\lim\limits_{x\to-\infty}f(x)=0$}.}{même démarche}
\explique{La droite d'équation $y=0$ est asymptote à $\Cf$ en $-\infty$ et en
$+\infty$.}{limite finie en l'infini}
\cas{b)}
\explique{On pose $u(x)=-x^{2}+2x$ : $u$ est dérivable sur $\R$ et $u'(x)=-2x+2$. Donc
$f=\e^{u}$ est dérivable sur $\R$ et, pour tout $x\in\R$ :}{forme $\e^{u}$}
\begin{calculs}
f'(x) &= u'(x)\e^{u(x)} &\qquad& \rem{$\left(\e^{u}\right)'=u'\e^{u}$}\\
f'(x) &= \res{(-2x+2)\e^{-x^{2}+2x}} &&
\end{calculs}
\cas{c)}
\explique{Pour tout $x\in\R$, $\e^{-x^{2}+2x}>0$, donc $f'(x)$ est du signe de
$-2x+2$, qui s'annule en $1$ : positif pour $x<1$, négatif pour $x>1$.}{une
exponentielle est toujours strictement positive}
\begin{calculs}
f(1) &= \e^{-1^{2}+2\times 1} = \e^{1} &\qquad& \rem{la valeur du maximum}
\end{calculs}
\begin{center}
\begin{tikzpicture}[scale=0.9]
  \tkzTabInit[lgt=2,espcl=2.6]{$x$/0.8,$f'(x)$/0.8,$f(x)$/1.6}{$-\infty$,$1$,$+\infty$}
  \tkzTabLine{,+,z,-,}
  \tkzTabVar{-/$0$,+/$\e$,-/$0$}
\end{tikzpicture}
\end{center}
\explique{$f$ est \res{strictement croissante sur $\intof{-\infty}{1}$, strictement
décroissante sur $\intfo{1}{+\infty}$}, de maximum $f(1)=\e$.}{on conclut par une
phrase}
\end{correction}

\etape{À vous de jouer}
Soit la fonction $g$ définie sur $\R$ par $g(x)=\e^{x^{2}-4x}$.
\begin{questions}[label=\alph*)]
  \item Déterminer les limites de $g$ en $-\infty$ et en $+\infty$.
  \item Calculer $g'(x)$ pour tout $x\in\R$.
  \item Dresser le tableau de variations de $g$.
\end{questions}
\reponse{a) $+\infty$ en $-\infty$ et en $+\infty$ ;\; b)~$g'(x)=(2x-4)\e^{x^{2}-4x}$ ;\;
c)~$g$ strictement décroissante sur $\intof{-\infty}{2}$, strictement croissante sur
$\intfo{2}{+\infty}$, de minimum $g(2)=\e^{-4}$.}
\end{methodebox}

\afaire{exercices 10 à 12, puis l\'entraînement, exercice 49}

\partiecours{Dérivée seconde}

La dérivée $f'$ est elle-même une fonction : on peut, à son tour, la dériver. Son
signe dira comment varie $f'$, c'est-à-dire comment « tourne » la courbe de $f$
(partie~V).

\begin{definitionbox}[Définition 2 --- Dérivée seconde]
Soit $f$ une fonction dérivable sur un intervalle $I$, dont la dérivée $f'$ est
dérivable sur $I$. On dit que $f$ est \textbf{deux fois dérivable} sur $I$ ; la dérivée
de $f'$ s'appelle la \textbf{dérivée seconde} de $f$, notée $f''$ :
\begin{center}
$f''(x)=\left(f'\right)'(x)$ pour tout $x\in I$.
\end{center}
\end{definitionbox}

\begin{methodebox}[Méthode 6 --- Calculer une dérivée seconde]
\etapeexemples
Calculer $f''(x)$ et $g''(x)$ pour tout $x\in\R$.
\begin{questions}[label=\alph*)]
  \item La fonction $f$ définie sur $\R$ par $f(x)=2x^{3}-9x^{2}+4$.
  \item La fonction $g$ définie sur $\R$ par $g(x)=(x+2)\e^{-x}$.
\end{questions}
\begin{correction}
\cas{a)}
\begin{calculs}
f'(x) &= 6x^{2}-18x &\qquad& \rem{on dérive terme à terme}\\
f''(x) &= \res{12x-18} &\qquad& \rem{on dérive $f'$}
\end{calculs}
\cas{b)}
\explique{$g$ est deux fois dérivable sur $\R$, comme produit de fonctions qui le sont.
Pour tout $x\in\R$ :}{$u(x)=x+2$ et $v(x)=\e^{-x}$}
\begin{calculs}
g'(x) &= 1\times\e^{-x}+(x+2)\times\left(-\e^{-x}\right) &\qquad& \rem{$u'v+uv'$}\\
g'(x) &= (1-x-2)\e^{-x} &\qquad& \rem{on factorise par $\e^{-x}$}\\
g'(x) &= (-x-1)\e^{-x} &\qquad& \rem{forme factorisée : on la dérive à son tour}\\
g''(x) &= -1\times\e^{-x}+(-x-1)\times\left(-\e^{-x}\right) &\qquad& \rem{$u'v+uv'$ avec $u(x)=-x-1$}\\
g''(x) &= (-1+x+1)\e^{-x} &&\\
g''(x) &= \res{x\e^{-x}} &&
\end{calculs}
\end{correction}

\etape{À vous de jouer}
Calculer $f''(x)$ et $g''(x)$ pour tout $x\in\R$.
\begin{questions}[label=\alph*)]
  \item La fonction $f$ définie sur $\R$ par $f(x)=x^{3}+3x^{2}-2x+1$.
  \item La fonction $g$ définie sur $\R$ par $g(x)=(x-1)\e^{-x}$.
\end{questions}
\reponse{a) $f''(x)=6x+6$ ;\; b)~$g''(x)=(x-3)\e^{-x}$.}
\end{methodebox}

\afaire{exercices 13 à 15}

\partiecours{Fonctions convexes, fonctions concaves}

Une \textbf{corde} d'une courbe est un segment qui relie deux de ses points. Certaines
courbes sont « tournées vers le haut », comme un bol : chaque corde passe au-dessus de
la courbe. D'autres sont « tournées vers le bas », comme une voûte.

\begin{definitionbox}[Définition 3 --- Fonction convexe, fonction concave]
\begin{minipage}[c]{0.55\linewidth}\raggedright
Soit $f$ une fonction définie sur un intervalle $I$.
\begin{itemize}
  \item $f$ est \textbf{convexe} sur $I$ si, pour tous points $A$ et $B$ de $\Cf$
        d'abscisses dans $I$, la courbe est située \textbf{en dessous} de la corde
        $[AB]$ entre $A$ et $B$.
  \item $f$ est \textbf{concave} sur $I$ si, pour tous points $A$ et $B$ de $\Cf$
        d'abscisses dans $I$, la courbe est située \textbf{au-dessus} de la corde
        $[AB]$ entre $A$ et $B$.
\end{itemize}
\end{minipage}\hfill
\begin{minipage}[c]{0.42\linewidth}\centering
\begin{tikzpicture}[x=0.42cm,y=0.42cm]
  \draw[-{Stealth[length=1.6mm]}] (0,0)--(6.3,0);
  \draw[-{Stealth[length=1.6mm]}] (0,0)--(0,4.3);
  \draw[coulCourbe,line width=1.1pt,domain=0.3:5.8,samples=60] plot (\x,{0.25*(\x-3)*(\x-3)+0.5});
  \draw[coulCourbeB,line width=0.9pt] (1,1.5)--(5.5,2.0625);
  \path[coulCourbeB] (1,1.5) node[croix]{} node[above left,font=\footnotesize]{$A$}
    (5.5,2.0625) node[croix]{} node[below right,font=\footnotesize]{$B$};
  \node[font=\footnotesize] at (3,-0.7) {convexe};
\end{tikzpicture}\hspace{0.4cm}%
\begin{tikzpicture}[x=0.42cm,y=0.42cm]
  \draw[-{Stealth[length=1.6mm]}] (0,0)--(6.3,0);
  \draw[-{Stealth[length=1.6mm]}] (0,0)--(0,4.3);
  \draw[coulCourbe,line width=1.1pt,domain=0.3:5.8,samples=60] plot (\x,{-0.25*(\x-3)*(\x-3)+3.5});
  \draw[coulCourbeB,line width=0.9pt] (1,2.5)--(5.5,1.9375);
  \path[coulCourbeB] (1,2.5) node[croix]{} node[below left,font=\footnotesize]{$A$}
    (5.5,1.9375) node[croix]{} node[below left,font=\footnotesize]{$B$};
  \node[font=\footnotesize] at (3,-0.7) {concave};
\end{tikzpicture}
\end{minipage}
\end{definitionbox}

\begin{proprietebox}[Théorème 2 --- Convexité d'une fonction deux fois dérivable (admis)]
Soit $f$ une fonction deux fois dérivable sur un intervalle $I$. Les quatre affirmations
suivantes sont équivalentes :
\begin{center}
$f$ est convexe sur $I$ \quad$\iff$\quad $\Cf$ est au-dessus de toutes ses tangentes sur
$I$\\[3pt]
$\iff$\quad $f'$ est croissante sur $I$ \quad$\iff$\quad $f''(x)\geq 0$ pour tout
$x\in I$.
\end{center}
De même : $f$ concave sur $I$ $\iff$ $\Cf$ en dessous de ses tangentes $\iff$ $f'$
décroissante sur $I$ $\iff$ $f''(x)\leq 0$ pour tout $x\in I$.

\smallskip
Fonctions de référence : $x\mapsto x^{2}$ et $x\mapsto\e^{x}$ sont convexes sur $\R$ ;
$x\mapsto x^{3}$ est concave sur $\intof{-\infty}{0}$ et convexe sur
$\intfo{0}{+\infty}$ ; $x\mapsto\frac{1}{x}$ est concave sur $\intoo{-\infty}{0}$ et
convexe sur $\intoo{0}{+\infty}$ ; $x\mapsto\sqrt{x}$ est concave sur
$\intfo{0}{+\infty}$.
\end{proprietebox}

\textbf{\color{coulMeth}Démonstration} (exigible) : si $f''(x)\geq 0$ pour tout $x\in I$,
alors $\Cf$ est au-dessus de ses tangentes.

\smallskip
Soit $a\in I$. La tangente $(T_{a})$ a pour équation $y=f'(a)(x-a)+f(a)$. On étudie
l'écart entre la courbe et la tangente : soit $g$ la fonction définie sur $I$ par
$g(x)=f(x)-\bigl[f'(a)(x-a)+f(a)\bigr]$. Elle est dérivable sur $I$ et
\begin{calculs}
g'(x) &= f'(x)-f'(a) &\qquad& \rem{$f'(a)(x-a)+f(a)$ est affine, de coefficient $f'(a)$}\\
g(a) &= f(a)-\bigl[0+f(a)\bigr] = 0 &\qquad& \rem{la tangente passe par le point d'abscisse $a$}
\end{calculs}
Comme $f''\geq 0$ sur $I$, $f'$ est croissante sur $I$. Donc, pour $x\leq a$,
$f'(x)\leq f'(a)$, soit $g'(x)\leq 0$ ; pour $x\geq a$, $g'(x)\geq 0$.

\begin{center}
\begin{tikzpicture}[scale=0.85]
  \tkzTabInit[lgt=2,espcl=2.4]{$x$/0.8,$g'(x)$/0.8,$g(x)$/1.4}{,$a$,}
  \tkzTabLine{,-,z,+,}
  \tkzTabVar{+/{},-/$0$,+/{}}
\end{tikzpicture}
\end{center}
Le minimum de $g$ sur $I$ est $g(a)=0$ : pour tout $x\in I$, $g(x)\geq 0$, soit
$f(x)\geq f'(a)(x-a)+f(a)$. La courbe $\Cf$ est au-dessus de $(T_{a})$, et cela pour
tout $a\in I$.

\begin{methodebox}[Méthode 7 --- Lire la convexité sur un graphique]
\etapeexemples
\begin{minipage}[c]{0.48\linewidth}\centering
\begin{tikzpicture}[x=0.5cm,y=0.5cm]
  \repere{-3}{-1}{3}{4}
  \gradx{-2,-1,1,2} \grady{1,2,3}
  \draw[coulCourbe,line width=1.1pt]
    (-2,1) .. controls (-1.67,2) and (-1.33,3) .. (-1,3)
           .. controls (-0.67,3) and (-0.33,2.33) .. (0,2)
           .. controls (0.33,1.67) and (0.67,1) .. (1,1)
           .. controls (1.33,1) and (1.67,2) .. (2,3);
  \node[coulCourbe,font=\footnotesize] at (2.5,3.6) {$\Cf$};
\end{tikzpicture}

{\small a) Courbe $\Cf$ d'une fonction $f$ définie sur $\intff{-2}{2}$.}
\end{minipage}\hfill
\begin{minipage}[c]{0.48\linewidth}\centering
\begin{tikzpicture}[x=0.5cm,y=0.5cm]
  \repere{-2}{-2}{4}{4}
  \gradx{-1,1,2,3} \grady{-1,1,2,3}
  \draw[coulCourbe,line width=1.1pt]
    (-1,3) .. controls (-0.67,1.67) and (-0.33,0.67) .. (0,0)
           .. controls (0.33,-0.67) and (0.67,-1) .. (1,-1)
           .. controls (1.33,-1) and (1.67,-0.67) .. (2,0)
           .. controls (2.33,0.67) and (2.67,1.67) .. (3,3);
  \node[coulCourbe,font=\footnotesize] at (3.6,3.4) {$\mathcal{C}_{g'}$};
\end{tikzpicture}

{\small b) Courbe $\mathcal{C}_{g'}$ de la \textbf{dérivée} $g'$ d'une fonction $g$
définie sur $\intff{-1}{3}$.}
\end{minipage}

\smallskip
Dans chaque cas, conjecturer les intervalles sur lesquels la fonction ($f$, puis $g$)
est convexe ou concave.
\begin{correction}
\cas{a)}\par\nopagebreak
\begin{minipage}[c]{0.58\linewidth}\raggedright
\explique{Sur $\intff{-2}{0}$, la courbe est en dessous de ses tangentes (par exemple
celle d'équation $y=3$) : \res{$f$ semble concave sur $\intff{-2}{0}$}.}{tournée vers le
bas}
\explique{Sur $\intff{0}{2}$, la courbe est au-dessus de ses tangentes (par exemple
celle d'équation $y=1$) : \res{$f$ semble convexe sur $\intff{0}{2}$}.}{tournée vers le
haut}
\end{minipage}\hfill
\begin{minipage}[c]{0.38\linewidth}\centering
\begin{tikzpicture}[x=0.45cm,y=0.45cm]
  \repere{-3}{-1}{3}{4}
  \gradx{-2,-1,1,2} \grady{1,2,3}
  \draw[coulCourbeB,dashed,line width=0.8pt] (-2.4,3)--(0.4,3) (-0.4,1)--(2.4,1);
  \draw[coulCourbe,line width=1.1pt]
    (-2,1) .. controls (-1.67,2) and (-1.33,3) .. (-1,3)
           .. controls (-0.67,3) and (-0.33,2.33) .. (0,2)
           .. controls (0.33,1.67) and (0.67,1) .. (1,1)
           .. controls (1.33,1) and (1.67,2) .. (2,3);
  \node[coulCourbe,font=\footnotesize] at (2.5,3.6) {$\Cf$};
\end{tikzpicture}
\end{minipage}
\cas{b)}
\explique{On lit le \textbf{sens de variation} de $g'$, pas la forme de sa courbe :
$g'$ est décroissante sur $\intff{-1}{1}$, puis croissante sur $\intff{1}{3}$.}{$g$
convexe $\iff$ $g'$ croissante}
\explique{Donc \res{$g$ est concave sur $\intff{-1}{1}$ et convexe sur
$\intff{1}{3}$}.}{théorème~2}
\end{correction}

\etape{À vous de jouer}
\begin{minipage}[c]{0.48\linewidth}\centering
\begin{tikzpicture}[x=0.5cm,y=0.5cm]
  \repere{-3}{-1}{3}{4}
  \gradx{-2,-1,1,2} \grady{1,2,3}
  \draw[coulCourbe,line width=1.1pt]
    (-2,3) .. controls (-1.67,2) and (-1.33,1) .. (-1,1)
           .. controls (-0.67,1) and (-0.33,1.67) .. (0,2)
           .. controls (0.33,2.33) and (0.67,3) .. (1,3)
           .. controls (1.33,3) and (1.67,2) .. (2,1);
  \node[coulCourbe,font=\footnotesize] at (2.5,3.6) {$\Cf$};
\end{tikzpicture}

{\small a) Courbe $\Cf$ d'une fonction $f$ définie sur $\intff{-2}{2}$.}
\end{minipage}\hfill
\begin{minipage}[c]{0.48\linewidth}\centering
\begin{tikzpicture}[x=0.5cm,y=0.5cm]
  \repere{-2}{-3}{4}{3}
  \gradx{-1,1,2,3} \grady{-2,-1,1,2}
  \draw[coulCourbe,line width=1.1pt]
    (-1,-2) .. controls (-0.67,-0.67) and (-0.33,0.33) .. (0,1)
            .. controls (0.33,1.67) and (0.67,2) .. (1,2)
            .. controls (1.33,2) and (1.67,1.67) .. (2,1)
            .. controls (2.33,0.33) and (2.67,-0.67) .. (3,-2);
  \node[coulCourbe,font=\footnotesize] at (3.6,-2.6) {$\mathcal{C}_{g'}$};
\end{tikzpicture}

{\small b) Courbe $\mathcal{C}_{g'}$ de la \textbf{dérivée} $g'$ d'une fonction $g$
définie sur $\intff{-1}{3}$.}
\end{minipage}

\smallskip
Dans chaque cas, conjecturer les intervalles sur lesquels la fonction ($f$, puis $g$)
est convexe ou concave.
\reponse{a) $f$ convexe sur $\intff{-2}{0}$, concave sur $\intff{0}{2}$ ;\; b)~$g$
convexe sur $\intff{-1}{1}$, concave sur $\intff{1}{3}$.}
\end{methodebox}

\begin{methodebox}[Méthode 8 --- Étudier la convexité par le signe de la dérivée seconde]
\etapeexemples
Étudier la convexité de chaque fonction sur $\R$.
\begin{questions}[label=\alph*)]
  \item La fonction $g$ définie sur $\R$ par $g(x)=(x+2)\e^{-x}$, dont la dérivée
        seconde est $g''(x)=x\e^{-x}$ (méthode~6).
  \item La fonction $h$ définie sur $\R$ par $h(x)=\e^{-x^{2}}$.
\end{questions}
\begin{correction}
\cas{a)}
\explique{Pour tout $x\in\R$, $\e^{-x}>0$, donc $g''(x)$ est du signe de $x$.}{on
étudie le signe de $g''$}
\begin{center}
\begin{tikzpicture}[scale=0.9]
  \tkzTabInit[lgt=3.4,espcl=2.8]{$x$/0.8,$g''(x)$/0.8,convexité de $g$/0.8}{$-\infty$,$0$,$+\infty$}
  \tkzTabLine{,-,z,+,}
  \tkzTabLine{,\text{concave},t,\text{convexe},}
\end{tikzpicture}
\end{center}
\explique{\res{$g$ est concave sur $\intof{-\infty}{0}$ et convexe sur
$\intfo{0}{+\infty}$}.}{théorème~2}
\cas{b)}
\explique{On pose $u(x)=-x^{2}$ : $u'(x)=-2x$. $h$ est deux fois dérivable sur $\R$ et,
pour tout $x\in\R$ :}{on calcule d'abord $h''$}
\begin{calculs}
h'(x) &= -2x\e^{-x^{2}} &\qquad& \rem{$\left(\e^{u}\right)'=u'\e^{u}$}\\
h''(x) &= -2\times\e^{-x^{2}}+(-2x)\times\left(-2x\e^{-x^{2}}\right) &\qquad& \rem{produit de $-2x$ par $\e^{-x^{2}}$}\\
h''(x) &= \left(4x^{2}-2\right)\e^{-x^{2}} &\qquad& \rem{on factorise par $\e^{-x^{2}}$}
\end{calculs}
\explique{$\e^{-x^{2}}>0$, donc $h''(x)$ est du signe du trinôme $4x^{2}-2$, de racines
$-\frac{\sqrt{2}}{2}$ et $\frac{\sqrt{2}}{2}$ ($x^{2}=\frac{1}{2}$), du signe de $a=4$ à
l'extérieur des racines.}{signe d'un trinôme (1re)}
\begin{center}
\begin{tikzpicture}[scale=0.9]
  \tkzTabInit[lgt=3.4,espcl=2.4]{$x$/0.8,$h''(x)$/0.8,convexité de $h$/0.8}{$-\infty$,$-\frac{\sqrt{2}}{2}$,$\frac{\sqrt{2}}{2}$,$+\infty$}
  \tkzTabLine{,+,z,-,z,+,}
  \tkzTabLine{,\text{convexe},t,\text{concave},t,\text{convexe},}
\end{tikzpicture}
\end{center}
\explique{\res{$h$ est convexe sur $\intof{-\infty}{-\frac{\sqrt{2}}{2}}$ et sur
$\intfo{\frac{\sqrt{2}}{2}}{+\infty}$, concave sur
$\intff{-\frac{\sqrt{2}}{2}}{\frac{\sqrt{2}}{2}}$}.}{}
\end{correction}

\etape{À vous de jouer}
Étudier la convexité de chaque fonction sur $\R$.
\begin{questions}[label=\alph*)]
  \item La fonction $g$ définie sur $\R$ par $g(x)=(x-1)\e^{-x}$, dont la dérivée
        seconde est $g''(x)=(x-3)\e^{-x}$.
  \item La fonction $h$ définie sur $\R$ par $h(x)=\e^{-2x^{2}}$.
\end{questions}
\reponse{a) $g$ concave sur $\intof{-\infty}{3}$, convexe sur $\intfo{3}{+\infty}$ ;\;
b)~$h''(x)=\left(16x^{2}-4\right)\e^{-2x^{2}}$ : $h$ convexe sur
$\intof{-\infty}{-\frac{1}{2}}$ et sur $\intfo{\frac{1}{2}}{+\infty}$, concave sur
$\intff{-\frac{1}{2}}{\frac{1}{2}}$.}
\end{methodebox}

\afaire{exercices 16 à 20}

\partiecours{Point d'inflexion}

Là où une fonction passe de concave à convexe (ou l'inverse), la courbe change de côté
par rapport à ses tangentes : elle traverse la tangente en ce point.

\begin{definitionbox}[Définition 4 --- Point d'inflexion]
\begin{minipage}[c]{0.58\linewidth}\raggedright
Soit $f$ une fonction dérivable sur un intervalle $I$, et $A$ un point de $\Cf$
d'abscisse $a\in I$. On dit que $A$ est un \textbf{point d'inflexion} de $\Cf$ si la
courbe \textbf{traverse} sa tangente en $A$ : en ce point, $f$ change de convexité.
\end{minipage}\hfill
\begin{minipage}[c]{0.38\linewidth}\centering
\begin{tikzpicture}[x=0.55cm,y=0.55cm]
  \draw[-{Stealth[length=1.6mm]}] (-3.2,0)--(3.4,0) node[below left]{$x$};
  \draw[-{Stealth[length=1.6mm]}] (0,-3.2)--(0,3.4) node[below left]{$y$};
  \draw[coulCourbeB,dashed,line width=0.8pt] (-3,-1.2)--(3,1.2);
  \draw[coulCourbe,line width=1.1pt,domain=-2.9:2.9,samples=60] plot (\x,{0.1*\x*\x*\x+0.4*\x});
  \path (0,0) node[croix]{} node[above left,font=\footnotesize]{$A$};
  \node[coulCourbe,font=\footnotesize] at (2.2,3.1) {$\Cf$};
  \node[coulCourbeB,font=\footnotesize] at (2.9,0.6) {$(T_{0})$};
\end{tikzpicture}
\end{minipage}
\end{definitionbox}

\begin{proprietebox}[Propriété 4 --- Point d'inflexion et dérivée seconde (admis)]
Soit $f$ une fonction deux fois dérivable sur un intervalle $I$ et $a\in I$. Si $f''$
\textbf{s'annule en changeant de signe} en $a$, alors le point de $\Cf$ de coordonnées
$\bigl(a\,;f(a)\bigr)$ est un point d'inflexion.

\smallskip
Il faut le changement de signe : pour $f(x)=x^{4}$, $f''(x)=12x^{2}$ s'annule en $0$
sans changer de signe, et $f$ est convexe sur $\R$ (pas d'inflexion).
\end{proprietebox}

\begin{methodebox}[Méthode 9 --- Déterminer un point d'inflexion]
\etapeexemples
Une entreprise fabrique des vélos électriques, jusqu'à $1\,000$ par mois. Le coût de
fabrication de $x$ centaines de vélos, en milliers d'euros, est modélisé par la fonction
$C$ définie sur $\intff{0}{10}$ par $C(x)=0{,}02x^{3}-0{,}3x^{2}+2x+10$.
\begin{questions}[label=\alph*)]
  \item Étudier la convexité de $C$ sur $\intff{0}{10}$.
  \item En déduire que la courbe de $C$ admet un point d'inflexion, dont on donnera les
        coordonnées.
  \item Interpréter ce résultat pour l'entreprise.
\end{questions}
\begin{correction}
\cas{a)}
\begin{calculs}
C'(x) &= 0{,}06x^{2}-0{,}6x+2 &\qquad& \rem{on dérive terme à terme}\\
C''(x) &= 0{,}12x-0{,}6 &&
\end{calculs}
\begin{calculs}
0{,}12x-0{,}6 &\geq 0 &\qquad& \rem{l'inéquation $C''(x)\geq 0$}\\
0{,}12x &\geq 0{,}6 &&\\
x &\geq 5 &\qquad& \rem{on divise par $0{,}12>0$}
\end{calculs}
\begin{center}
\begin{tikzpicture}[scale=0.9]
  \tkzTabInit[lgt=3.4,espcl=2.8]{$x$/0.8,$C''(x)$/0.8,convexité de $C$/0.8}{$0$,$5$,$10$}
  \tkzTabLine{,-,z,+,}
  \tkzTabLine{,\text{concave},t,\text{convexe},}
\end{tikzpicture}
\end{center}
\explique{\res{$C$ est concave sur $\intff{0}{5}$ et convexe sur $\intff{5}{10}$}.}{théorème~2}
\cas{b)}
\explique{$C''$ s'annule en changeant de signe en $5$ : la courbe admet un point
d'inflexion d'abscisse $5$.}{propriété~4}
\begin{calculs}
C(5) &= 0{,}02\times 125-0{,}3\times 25+2\times 5+10 &&\\
C(5) &= 2{,}5-7{,}5+10+10 = 15 &&
\end{calculs}
\explique{Le point d'inflexion a pour coordonnées \res{$(5\,;15)$}.}{}
\cas{c)}
\explique{$C''$ est négative puis positive : $C'$ est décroissante sur $\intff{0}{5}$,
croissante sur $\intff{5}{10}$, de minimum $C'(5)=1{,}5-3+2=0{,}5>0$. Donc $C$ est
croissante.}{$C''$ donne les variations de $C'$}
\explique{\res{Jusqu'à $500$ vélos, le coût augmente de moins en moins vite ; au-delà,
il augmente de plus en plus vite}.}{on répond dans le contexte}
\end{correction}

\etape{À vous de jouer}
Une entreprise fabrique des panneaux solaires, jusqu'à $2\,000$ par mois. Le coût de
fabrication de $x$ centaines de panneaux, en milliers d'euros, est modélisé par la
fonction $C$ définie sur $\intff{0}{20}$ par $C(x)=0{,}01x^{3}-0{,}3x^{2}+4x+20$.
\begin{questions}[label=\alph*)]
  \item Étudier la convexité de $C$ sur $\intff{0}{20}$.
  \item En déduire que la courbe de $C$ admet un point d'inflexion, dont on donnera les
        coordonnées.
  \item Interpréter ce résultat pour l'entreprise.
\end{questions}
\reponse{a) $C$ concave sur $\intff{0}{10}$, convexe sur $\intff{10}{20}$ ;\; b)~$(10\,;40)$ ;\;
c)~jusqu'à $1\,000$ panneaux, le coût augmente de moins en moins vite, puis de plus en
plus vite.}
\end{methodebox}

\afaire{exercices 21 à 23}

\partiecours{Convexité et inégalités}

Une fonction convexe a sa courbe au-dessus de ses tangentes (théorème~2) : en
comparant $f(x)$ à l'équation d'une tangente, on obtient une inégalité valable sur tout
un intervalle, sans étudier de nouvelle fonction.

\begin{methodebox}[Méthode 10 --- Démontrer une inégalité avec la convexité]
\etapeexemples
Soit la fonction $f$ définie sur $\R$ par $f(x)=\e^{x}$.
\begin{questions}[label=\alph*)]
  \item Étudier la convexité de $f$ sur $\R$.
  \item Déterminer une équation de la tangente $(T_{0})$ à $\Cf$ au point d'abscisse $0$.
  \item En déduire que, pour tout $x\in\R$, $\e^{x}\geq x+1$.
\end{questions}
\begin{correction}
\cas{a)}
\explique{Pour tout $x\in\R$, $f'(x)=\e^{x}$ et $f''(x)=\e^{x}>0$ : \res{$f$ est
convexe sur $\R$}.}{théorème~2}
\cas{b)}
\begin{calculs}
f(0) &= \e^{0} = 1 &&\\
f'(0) &= \e^{0} = 1 &&\\
y &= 1\times(x-0)+1 &\qquad& \rem{$y=f'(0)(x-0)+f(0)$}\\
y &= \res{x+1} &&
\end{calculs}
\cas{c)}
\explique{$f$ est convexe sur $\R$, donc sa courbe est au-dessus de toutes ses
tangentes, en particulier de $(T_{0})$, d'équation $y=x+1$.}{la rédaction des corrigés
de bac}
\explique{Donc, \res{pour tout $x\in\R$, $\e^{x}\geq x+1$}.}{l'ordonnée sur la courbe
dépasse celle sur la tangente}
\end{correction}

\etape{À vous de jouer}
Soit la fonction $f$ définie sur $\intoo{0}{+\infty}$ par $f(x)=\sqrt{x}$.
\begin{questions}[label=\alph*)]
  \item Étudier la convexité de $f$ sur $\intoo{0}{+\infty}$.
  \item Déterminer une équation de la tangente $(T_{1})$ à $\Cf$ au point d'abscisse $1$.
  \item En déduire que, pour tout $x>0$, $\sqrt{x}\leq\dfrac{x+1}{2}$.
\end{questions}
\reponse{a) $f''(x)=-\dfrac{1}{4x\sqrt{x}}<0$ : $f$ concave ;\;
b)~$(T_{1})\colon y=\frac{1}{2}x+\frac{1}{2}$ ;\; c)~courbe en dessous de $(T_{1})$.}
\end{methodebox}

\afaire{exercices 24 à 26}

\partiecours{Continuité}

Un réservoir se remplit d'eau : son volume passe de $0$ à $100$ litres sans « sauter »
de valeur. Le prix d'un colis, lui, saute d'un tarif au suivant quand la masse dépasse
$2$~kg. La continuité distingue ces deux comportements.

\smallskip
Exemple : soit la fonction $f$ définie sur $\R$ par $f(x)=x^{2}$. On a $f(1{,}9)=3{,}61$ ;
$f(1{,}99)=3{,}960\,1$ et $f(1{,}999)=3{,}996\,001$ : quand $x$ s'approche de $2$,
$f(x)$ s'approche de $4$, qui est $f(2)$. Pas de saut en $2$.

\begin{definitionbox}[Définition 5 --- Fonction continue]
\begin{minipage}[c]{0.55\linewidth}\raggedright
\textbf{Intuitivement :} une fonction est continue sur un intervalle si sa courbe s'y
trace \textbf{sans lever le crayon}.

\smallskip
\textbf{Définition :} soit $f$ une fonction définie sur un intervalle $I$ et $a\in I$.
\begin{itemize}
  \item $f$ est \textbf{continue en $a$} si $\lim\limits_{x\to a}f(x)=f(a)$ ;
  \item $f$ est \textbf{continue sur $I$} si elle est continue en tout nombre réel de
        $I$.
\end{itemize}
\end{minipage}\hfill
\begin{minipage}[c]{0.42\linewidth}\centering
\begin{tikzpicture}[x=0.4cm,y=0.4cm]
  \draw[-{Stealth[length=1.6mm]}] (0,0)--(5.4,0);
  \draw[-{Stealth[length=1.6mm]}] (0,0)--(0,4.4);
  \draw[coulCourbe,line width=1.1pt,domain=0.3:5,samples=40] plot (\x,{0.12*\x*\x+0.4*\x+0.5});
  \draw[dotted] (2.5,0)--(2.5,2.25);
  \path (2.5,2.25) node[croix]{};
  \node[below,font=\footnotesize] at (2.5,0) {$a$};
  \node[font=\footnotesize,align=center] at (2.7,-1.4) {continue en $a$};
\end{tikzpicture}\hspace{0.3cm}%
\begin{tikzpicture}[x=0.4cm,y=0.4cm]
  \draw[-{Stealth[length=1.6mm]}] (0,0)--(5.4,0);
  \draw[-{Stealth[length=1.6mm]}] (0,0)--(0,4.4);
  \draw[coulCourbe,line width=1.1pt] (0.3,0.6)--(2.5,1.5);
  \draw[coulCourbe,line width=1.1pt] (2.5,3)--(5,3.8);
  \draw[dotted] (2.5,0)--(2.5,3);
  \path (2.5,3) node[croix]{};
  \node[below,font=\footnotesize] at (2.5,0) {$a$};
  \node[font=\footnotesize,align=center] at (2.7,-1.4) {pas continue en $a$};
\end{tikzpicture}
\end{minipage}
\end{definitionbox}

\begin{proprietebox}[Propriété 5 --- Fonctions continues (admis)]
\begin{itemize}
  \item Les fonctions polynômes, la fonction exponentielle et la fonction valeur absolue
        sont continues sur $\R$ ; la fonction inverse est continue sur
        $\intoo{-\infty}{0}$ et sur $\intoo{0}{+\infty}$ ; la fonction racine carrée est
        continue sur $\intfo{0}{+\infty}$.
  \item Une somme, un produit, un quotient (dont le dénominateur ne s'annule pas) ou une
        composée de fonctions continues est continue.
  \item \textbf{Une fonction dérivable sur un intervalle $I$ est continue sur $I$.} La
        réciproque est fausse : la fonction racine carrée est continue en $0$, mais pas
        dérivable en $0$ (tangente verticale).
\end{itemize}
Convention : dans un tableau de variations, une flèche oblique signifie que la fonction
est continue et strictement monotone sur l'intervalle.
\end{proprietebox}

\begin{proprietebox}[Propriété 6 --- Image d'une suite convergente (admis)]
Soit $(u_{n})$ une suite qui converge vers un nombre réel $\ell$, et $f$ une fonction
continue en $\ell$. Alors la suite $\bigl(f(u_{n})\bigr)$ converge vers $f(\ell)$.

\smallskip
Par exemple, la suite $(u_{n})$ définie pour tout $n\in\Ns$ par $u_{n}=2+\frac{1}{n}$
converge vers $2$, et la fonction exponentielle est continue en $2$ : la suite
$\left(\e^{u_{n}}\right)$ converge vers $\e^{2}$. C'est cette propriété qui fonde le
théorème du point fixe du chapitre~2 : si $u_{n+1}=f(u_{n})$ et $u_{n}\to\ell$, avec $f$
continue, alors $f(u_{n})\to f(\ell)$ et $u_{n+1}\to\ell$, d'où $f(\ell)=\ell$.
\end{proprietebox}

\begin{methodebox}[Méthode 11 --- Étudier la continuité d'une fonction définie par morceaux]
\etapeexemples
\begin{minipage}[c]{0.58\linewidth}\raggedright
Soit la fonction $f$ définie sur $\R$ par
\begin{center}
$f(x)=\begin{cases}x+1 & \text{si } x<1\\ x^{2}-4x+5 & \text{si } 1\leq x<3\\
x-4 & \text{si } x\geq 3\end{cases}$
\end{center}
La fonction $f$ est-elle continue en $1$ ? en $3$ ? sur $\R$ ?
\end{minipage}\hfill
\begin{minipage}[c]{0.38\linewidth}\centering
\begin{tikzpicture}[x=0.45cm,y=0.45cm]
  \repere{-2}{-2}{5}{3}
  \gradx{-1,1,2,3,4} \grady{-1,1,2}
  \draw[coulCourbe,line width=1.1pt] (-2,-1)--(1,2);
  \draw[coulCourbe,line width=1.1pt]
    (1,2) .. controls (1.33,1.33) and (1.67,1) .. (2,1)
          .. controls (2.33,1) and (2.67,1.33) .. (3,2);
  \draw[coulCourbe,line width=1.1pt] (3,-1)--(5,1);
  \path[coulCourbe] (3,-1) node[croix]{};
  \node[coulCourbe,font=\footnotesize] at (-1.2,1.5) {$\Cf$};
\end{tikzpicture}
\end{minipage}
\begin{correction}
\cas{Résolution}\par\nopagebreak
\explique{Sur $\intoo{-\infty}{1}$, sur $\intoo{1}{3}$ et sur $\intoo{3}{+\infty}$,
$f(x)$ est donnée par un polynôme : $f$ est continue sur chacun de ces
intervalles.}{propriété~5 ; reste à étudier les « raccords »}
\explique{En $1$, $\lim\limits_{x\to 1^{-}}f(x)=\lim\limits_{x\to 1^{-}}(x+1)=2$ et
$\lim\limits_{x\to 1^{+}}f(x)=\lim\limits_{x\to 1^{+}}\left(x^{2}-4x+5\right)=2$ ; or
$f(1)=1-4+5=2$. Donc $\lim\limits_{x\to 1}f(x)=f(1)$ : \res{$f$ est continue en
$1$}.}{limite à gauche, limite à droite, valeur}
\explique{En $3$, $\lim\limits_{x\to 3^{-}}f(x)=9-12+5=2$, mais $f(3)=3-4=-1$. Donc
$f(x)$ ne tend pas vers $f(3)$ : \res{$f$ n'est pas continue en $3$}.}{la courbe
« saute » en $3$}
\explique{\res{$f$ n'est pas continue sur $\R$} ; elle est continue sur
$\intoo{-\infty}{3}$ et sur $\intfo{3}{+\infty}$.}{}
\end{correction}

\etape{À vous de jouer}
Soit la fonction $g$ définie sur $\R$ par
\begin{center}
$g(x)=\begin{cases}-x+3 & \text{si } x<1\\ x^{2}-2x+3 & \text{si } 1\leq x<3\\
x+1 & \text{si } x\geq 3\end{cases}$
\end{center}
La fonction $g$ est-elle continue en $1$ ? en $3$ ? sur $\R$ ?
\reponse{continue en $1$ ; pas continue en $3$ ; pas continue sur $\R$.}
\end{methodebox}

\afaire{exercices 27 à 30}

\partiecours{Théorème des valeurs intermédiaires}

Une randonneuse part d'un refuge à $1\,200$~m d'altitude et arrive à un col à
$2\,000$~m : à un moment au moins, elle est passée par l'altitude $1\,500$~m, puisque son
altitude varie sans sauter de valeur. Elle a pu y passer plusieurs fois, si le
sentier monte et descend.

\begin{proprietebox}[Théorème 3 --- Théorème des valeurs intermédiaires (admis)]
\begin{minipage}[c]{0.55\linewidth}\raggedright
Soit $f$ une fonction \textbf{continue} sur un intervalle $\intff{a}{b}$. Pour tout
nombre réel $k$ compris entre $f(a)$ et $f(b)$, il existe \textbf{au moins}
un nombre réel $c\in\intff{a}{b}$ tel que $f(c)=k$.

\smallskip
Autrement dit, l'équation $f(x)=k$ admet au moins une solution dans $\intff{a}{b}$. Sur
le schéma, elle en a trois.
\end{minipage}\hfill
\begin{minipage}[c]{0.42\linewidth}\centering
\begin{tikzpicture}[x=1cm,y=0.8cm]
  \draw[-{Stealth[length=1.6mm]}] (0,0)--(4.5,0) node[below left]{$x$};
  \draw[-{Stealth[length=1.6mm]}] (0,0)--(0,4.2) node[below left]{$y$};
  \draw[dashed,black!60] (0,2.15)--(4.2,2.15);
  \draw[dotted] (0,1.6875)--(0.5,1.6875) (0.5,0)--(0.5,1.6875)
                (0,3.7)--(4,3.7) (4,0)--(4,3.7);
  \draw[coulCourbe,line width=1.1pt,domain=0.5:4,samples=80] plot (\x,{0.3*\x*\x*\x-1.8*\x*\x+3.2*\x+0.5});
  \foreach \c/\n in {0.91/1,1.87/2,3.22/3}{
    \draw[dotted,red] (\c,0)--(\c,2.15);
    \path[red] (\c,0) node[croix]{} node[below,font=\footnotesize]{$c_{\n}$};}
  \node[below,font=\footnotesize] at (0.5,0) {$a$};
  \node[below,font=\footnotesize] at (4,0) {$b$};
  \node[left,font=\footnotesize] at (0,1.6875) {$f(a)$};
  \node[left,font=\footnotesize] at (0,2.15) {$k$};
  \node[left,font=\footnotesize] at (0,3.7) {$f(b)$};
\end{tikzpicture}
\end{minipage}
\end{proprietebox}

\begin{proprietebox}[Théorème 4 --- Corollaire : fonction continue strictement monotone (admis)]
Soit $f$ une fonction \textbf{continue} et \textbf{strictement monotone} sur
$\intff{a}{b}$. Pour tout nombre réel $k$ compris entre $f(a)$ et $f(b)$, l'équation
$f(x)=k$ admet une \textbf{unique} solution dans $\intff{a}{b}$.

\smallskip
Ce corollaire reste vrai sur un intervalle ouvert ou semi-ouvert, comme
$\intfo{a}{+\infty}$ ou $\intoo{-\infty}{b}$ : on remplace $f(a)$ ou $f(b)$ par la
limite de $f$ à la borne ouverte. Cas fréquent : $k=0$, quand $f$ change de signe.
\end{proprietebox}

\begin{methodebox}[Méthode 12 --- Montrer qu'une équation a une unique solution, l'encadrer]
\etapeexemples
Soit la fonction $f$ définie sur $\R$ par $f(x)=x^{3}-3x^{2}+1$.
\begin{questions}[label=\alph*)]
  \item Montrer que l'équation $f(x)=0$ admet une unique solution $\alpha$ dans
        $\intff{2}{4}$.
  \item À l'aide de la calculatrice, donner un encadrement de $\alpha$ d'amplitude
        $10^{-2}$.
\end{questions}
\begin{correction}
\cas{a)}
\explique{$f$ est une fonction polynôme : elle est dérivable, donc continue, sur
$\intff{2}{4}$.}{première hypothèse : la continuité}
\begin{calculs}
f'(x) &= 3x^{2}-6x &&\\
f'(x) &= 3x(x-2) &\qquad& \rem{on factorise pour lire le signe}
\end{calculs}
\explique{Pour tout $x\in\intof{2}{4}$, $3x>0$ et $x-2>0$, donc $f'(x)>0$ ; $f'$ ne
s'annule qu'en $2$. Donc $f$ est strictement croissante sur $\intff{2}{4}$.}{deuxième
hypothèse : la stricte monotonie}
\begin{calculs}
f(2) &= 8-12+1 = -3 &&\\
f(4) &= 64-48+1 = 17 &&
\end{calculs}
\explique{$f$ est continue et strictement croissante sur $\intff{2}{4}$, et
$0\in\intff{-3}{17}$, c'est-à-dire $0$ est compris entre $f(2)$ et $f(4)$. D'après le
corollaire du théorème des valeurs intermédiaires, \res{l'équation $f(x)=0$ admet une
unique solution $\alpha$ dans $\intff{2}{4}$}.}{troisième point : $k$ entre $f(a)$ et
$f(b)$ ; puis on cite le théorème}
\cas{b)}
\explique{Avec le tableau de valeurs de la calculatrice, au pas de $0{,}1$, puis de
$0{,}01$ :}{on cherche où $f$ change de signe}
\begin{calculs}
f(2{,}8) &\approx -0{,}57 &\qquad& \rem{$f(2{,}8)<0$}\\
f(2{,}9) &\approx 0{,}16 &\qquad& \rem{$f(2{,}9)>0$ : $2{,}8<\alpha<2{,}9$}\\
f(2{,}87) &\approx -0{,}07 &&\\
f(2{,}88) &\approx 0{,}005 &&
\end{calculs}
\explique{$f(2{,}87)<0<f(2{,}88)$ et $f$ est strictement croissante : \res{$2{,}87<\alpha<2{,}88$}.}{}
\end{correction}

\etape{À vous de jouer}
Soit la fonction $g$ définie sur $\R$ par $g(x)=x^{3}-6x^{2}+5$.
\begin{questions}[label=\alph*)]
  \item Montrer que l'équation $g(x)=0$ admet une unique solution $\alpha$ dans
        $\intff{4}{6}$.
  \item À l'aide de la calculatrice, donner un encadrement de $\alpha$ d'amplitude
        $10^{-2}$.
\end{questions}
\reponse{a) $g'(x)=3x(x-4)>0$ sur $\intof{4}{6}$, $g(4)=-27$ et $g(6)=5$ ;\;
b)~$5{,}85<\alpha<5{,}86$.}
\end{methodebox}

\begin{methodebox}[Méthode 13 --- Déterminer le nombre de solutions d'une équation]
\etapeexemples
\begin{minipage}[c]{0.5\linewidth}\raggedright
Soit la fonction $f$ définie sur $\R$ par $f(x)=(x+1)\e^{-x}$. On admet son tableau de
variations ci-contre. Déterminer le nombre de solutions sur $\R$ de chaque équation.
\begin{questions}[label=\alph*)]
  \item $f(x)=\frac{1}{2}$
  \item $f(x)=-1$
\end{questions}
\end{minipage}\hfill
\begin{minipage}[c]{0.46\linewidth}\centering
\begin{tikzpicture}[scale=0.85]
  \tkzTabInit[lgt=1.6,espcl=2.2]{$x$/0.8,$f(x)$/1.6}{$-\infty$,$0$,$+\infty$}
  \tkzTabVar{-/$-\infty$,+/$1$,-/$0$}
\end{tikzpicture}
\end{minipage}
\begin{correction}
\cas{a)}
\explique{Sur $\intof{-\infty}{0}$, $f$ est continue (car dérivable) et strictement
croissante ; de plus $\lim\limits_{x\to-\infty}f(x)=-\infty$ et $f(0)=1$. Or
$\frac{1}{2}\in\intof{-\infty}{1}$ : d'après le corollaire du théorème des valeurs
intermédiaires, l'équation $f(x)=\frac{1}{2}$ admet une unique solution dans
$\intof{-\infty}{0}$.}{on découpe selon les flèches du tableau}
\explique{Sur $\intfo{0}{+\infty}$, $f$ est continue et strictement décroissante, de
$f(0)=1$ vers la limite $0$ : elle prend ses valeurs dans $\intof{0}{1}$. Or
$\frac{1}{2}\in\intof{0}{1}$ : l'équation admet une unique solution dans
$\intfo{0}{+\infty}$.}{la limite $0$ n'est pas atteinte}
\explique{Donc l'équation $f(x)=\frac{1}{2}$ admet \res{exactement deux solutions} sur
$\R$.}{on additionne}
\cas{b)}
\explique{Sur $\intof{-\infty}{0}$, $f$ est continue, strictement croissante, et
$-1\in\intof{-\infty}{1}$ : l'équation $f(x)=-1$ admet une unique solution dans cet
intervalle.}{même rédaction}
\explique{Sur $\intfo{0}{+\infty}$, $f(x)\in\intof{0}{1}$, donc $f(x)>0>-1$ : aucune
solution.}{$-1$ n'est pas une valeur prise}
\explique{Donc l'équation $f(x)=-1$ admet \res{une unique solution} sur $\R$.}{}
\end{correction}

\etape{À vous de jouer}
\begin{minipage}[c]{0.5\linewidth}\raggedright
Soit la fonction $g$ définie sur $\R$ par $g(x)=(1-x)\e^{x}$. On admet son tableau de
variations ci-contre. Déterminer le nombre de solutions sur $\R$ de chaque équation.
\begin{questions}[label=\alph*)]
  \item $g(x)=\frac{1}{2}$
  \item $g(x)=-1$
\end{questions}
\end{minipage}\hfill
\begin{minipage}[c]{0.46\linewidth}\centering
\begin{tikzpicture}[scale=0.85]
  \tkzTabInit[lgt=1.6,espcl=2.2]{$x$/0.8,$g(x)$/1.6}{$-\infty$,$0$,$+\infty$}
  \tkzTabVar{-/$0$,+/$1$,-/$-\infty$}
\end{tikzpicture}
\end{minipage}
\reponse{a) deux solutions ;\; b)~une unique solution.}
\end{methodebox}

\afaire{exercices 31 à 36}

\partiecours{Méthode de dichotomie}

Le théorème des valeurs intermédiaires dit qu'une solution existe, sans la donner. La
\textbf{dichotomie} (« couper en deux ») l'encadre de plus en plus finement : si $f$ est
continue et change de signe sur $\intff{a}{b}$, on calcule $f$ au milieu
$m=\frac{a+b}{2}$, et on garde la moitié où $f$ change encore de signe. À chaque étape,
l'amplitude de l'encadrement est divisée par $2$.

\begin{remarquebox}[Algorithme de dichotomie en Python]
\begin{minipage}[t]{0.46\linewidth}\vspace{0pt}
\begin{python}
def f(x):
    return x**3 - 3*x**2 + 1

def dichotomie(a, b, e):
    while b - a > e:
        m = (a + b) / 2
        if f(a) * f(m) <= 0:
            b = m
        else:
            a = m
    return a, b
\end{python}
\end{minipage}\hfill
\begin{minipage}[t]{0.5\linewidth}\vspace{0pt}\raggedright
La boucle tourne tant que l'amplitude $b-a$ dépasse la précision \texttt{e}.

\smallskip
Si $f(a)\times f(m)\leq 0$, $f(a)$ et $f(m)$ sont de signes contraires (ou l'un est
nul) : la solution est dans $\intff{a}{m}$, on remplace $b$ par $m$. Sinon, elle est
dans $\intff{m}{b}$, on remplace $a$ par $m$.

\smallskip
La fonction renvoie les bornes d'un encadrement de la solution.
\end{minipage}
\end{remarquebox}

\begin{methodebox}[Méthode 14 --- Encadrer une solution par dichotomie]
\etapeexemples
On reprend la fonction $f$ définie sur $\R$ par $f(x)=x^{3}-3x^{2}+1$ et le programme
ci-dessus. L'équation $f(x)=0$ a une unique solution $\alpha$ dans $\intff{2}{4}$, et
$f(2)=-3$.
\begin{questions}[label=\alph*)]
  \item Faire tourner à la main l'appel \texttt{dichotomie(2, 4, 0.25)} : donner les
        valeurs successives de $a$, $b$ et $m$, puis ce que renvoie la fonction.
  \item Combien de passages dans la boucle faut-il pour obtenir un encadrement
        d'amplitude au plus $10^{-2}$ ?
\end{questions}
\begin{correction}
\cas{a)}
\explique{Ici $f(a)<0$ à chaque étape : le test $f(a)\times f(m)\leq 0$ revient à
$f(m)\geq 0$.}{$a$ ne prend que des valeurs où $f<0$}
\begin{center}
\renewcommand{\arraystretch}{1.3}\setlength{\tabcolsep}{10pt}
\begin{tabular}{|c|c|c|c|c|}
\hline
passage & $m$ & $f(m)$ & nouvel intervalle & $b-a$\\ \hline
$1$ & $3$ & $1>0$ & $\intff{2}{3}$ & $1$\\ \hline
$2$ & $2{,}5$ & $-2{,}125<0$ & $\intff{2{,}5}{3}$ & $0{,}5$\\ \hline
$3$ & $2{,}75$ & $\approx -0{,}89<0$ & $\intff{2{,}75}{3}$ & $0{,}25$\\ \hline
\end{tabular}
\end{center}
\explique{$b-a=0{,}25$ ne dépasse plus $0{,}25$ : la boucle s'arrête et la fonction
renvoie $a=2{,}75$ et $b=3$, soit \res{$2{,}75\leq\alpha\leq 3$}.}{la condition du
\texttt{while} est fausse}
\cas{b)}
\explique{Au départ, $b-a=2$ ; après $n$ passages, l'amplitude vaut $\frac{2}{2^{n}}$.
On cherche le plus petit $n\in\N$ tel que :}{l'amplitude est divisée par $2$ à chaque
passage}
\begin{calculs}
\frac{2}{2^{n}} &\leq 0{,}01 &&\\
2 &\leq 0{,}01\times 2^{n} &\qquad& \rem{on multiplie par $2^{n}>0$}\\
2^{n} &\geq 200 &\qquad& \rem{on divise par $0{,}01>0$}
\end{calculs}
\explique{$2^{7}=128<200$ et $2^{8}=256\geq 200$ : il faut \res{$8$ passages}.}{}
\end{correction}

\etape{À vous de jouer}
On remplace la première fonction du programme par la fonction $g$ définie sur $\R$ par
$g(x)=x^{3}-6x^{2}+5$ (\texttt{return x**3 - 6*x**2 + 5}). L'équation $g(x)=0$ a une
unique solution $\alpha$ dans $\intff{4}{6}$, et $g(4)=-27$.
\begin{questions}[label=\alph*)]
  \item Faire tourner à la main l'appel \texttt{dichotomie(4, 6, 0.25)} : donner les
        valeurs successives de $a$, $b$ et $m$, puis ce que renvoie la fonction.
  \item Combien de passages dans la boucle faut-il pour obtenir un encadrement
        d'amplitude au plus $10^{-3}$ ?
\end{questions}
\reponse{a) $m=5$ ; $5{,}5$ ; $5{,}75$ : la fonction renvoie $5{,}75$ et $6$ ;\;
b)~$11$ passages.}
\end{methodebox}

\afaire{exercices 37 à 39, puis les exercices type bac 40 à 42}

\end{document}
